% Options for packages loaded elsewhere
\PassOptionsToPackage{unicode}{hyperref}
\PassOptionsToPackage{hyphens}{url}
\documentclass[
  11pt,
]{article}
\usepackage{xcolor}
\usepackage[margin=25mm]{geometry}
\usepackage{amsmath,amssymb}
\setcounter{secnumdepth}{-\maxdimen} % remove section numbering
\usepackage{iftex}
\ifPDFTeX
  \usepackage[T1]{fontenc}
  \usepackage[utf8]{inputenc}
  \usepackage{textcomp} % provide euro and other symbols
\else % if luatex or xetex
  \usepackage{unicode-math} % this also loads fontspec
  \defaultfontfeatures{Scale=MatchLowercase}
  \defaultfontfeatures[\rmfamily]{Ligatures=TeX,Scale=1}
\fi
\usepackage{lmodern}
\ifPDFTeX\else
  % xetex/luatex font selection
\fi
% Use upquote if available, for straight quotes in verbatim environments
\IfFileExists{upquote.sty}{\usepackage{upquote}}{}
\IfFileExists{microtype.sty}{% use microtype if available
  \usepackage[]{microtype}
  \UseMicrotypeSet[protrusion]{basicmath} % disable protrusion for tt fonts
}{}
\makeatletter
\@ifundefined{KOMAClassName}{% if non-KOMA class
  \IfFileExists{parskip.sty}{%
    \usepackage{parskip}
  }{% else
    \setlength{\parindent}{0pt}
    \setlength{\parskip}{6pt plus 2pt minus 1pt}}
}{% if KOMA class
  \KOMAoptions{parskip=half}}
\makeatother
\setlength{\emergencystretch}{3em} % prevent overfull lines
\providecommand{\tightlist}{%
  \setlength{\itemsep}{0pt}\setlength{\parskip}{0pt}}
\usepackage{bookmark}
\IfFileExists{xurl.sty}{\usepackage{xurl}}{} % add URL line breaks if available
\urlstyle{same}
\hypersetup{
  hidelinks,
  pdfcreator={LaTeX via pandoc}}

\author{}
\date{}

\begin{document}

\section{Equivariant KK-theory and the Green--Julg
theorem}\label{equivariant-kk-theory-and-the-greenjulg-theorem}

\emph{Written by GPT-6.1 Sol (OpenAI). Public domain (CC0).}

A symmetry can leave two identical ordinary Fredholm indices with
different equivariant labels: the finite-dimensional kernel may carry
the trivial representation or a sign representation. We will follow
those labels from a two-element action into invariant operator matrices,
then through stabilization to the Green--Julg isomorphism. The same
matrix calculation also exposes the Haar normalization that a
crossed-product formula must preserve.

A group action adds two requirements to a Kasparov cycle. Its
representation must respect the action, and its operator must respect it
up to compact errors after the source algebra acts. For a compact group,
averaging removes the latter errors, while invariant connections and
partitions construct the product. For a discrete group, individual
expressions \(aU_g\) already belong to the full crossed product, leading
to a different cycle comparison. The analytic order below reflects these
prerequisites: the compact homotopy comparison in Theorem 2.3 is needed
in the Green--Julg proof, and the product is needed for
representation-ring and induction laws.

We retain the graded tensor conventions of \emph{Graded C*-algebras,
Clifford algebras and graded Hilbert modules}. We use the
non-equivariant connection calculus and product proofs in
\emph{Connections and the existence of the Kasparov product}, Sections
1--5, and the associativity and homotopy proofs in \emph{Homotopy,
associativity, the index pairing and KK-equivalence}, Sections 1--5. The
integrated covariant representation and regular-module constructions are
Theorem 2.3 of \emph{C*-dynamical systems and full crossed products} and
Proposition 2.1 and the absorption proof of \emph{Reduced crossed
products and Fell's absorption principle}. We prove the additional
equivariant steps below.

Throughout, groups are second countable locally compact Hausdorff
groups, algebras are separable, and Hilbert modules are countably
generated. Actions preserve grading. Haar measure on a compact group is
normalized to have mass one. A crossed product without a subscript is
the \textbf{full} crossed product. For graded coefficients, \(K_i(D)\)
means the graded group \(KK^i(\mathbb C,D)\); it agrees with ordinary
K-theory when the grading is trivial.

For compact groups, the Haar probability measure used below is proved in
\href{supporting/representations-of-compact-groups/compact-foundations.html\#cpt-f-005}{Foundations
for compact-group averaging and coefficient approximation, CPT-F-005}.
Its
\href{supporting/representations-of-compact-groups/compact-foundations.html\#cpt-f-010}{CPT-F-010}
constructs continuous vector integrals, strict operator averages and
source-localized compact integrals on arbitrary Hilbert modules. For the
noncompact case,
\href{supporting/noncompact-foundations/noncompact-foundations.html}{Noncompact
foundations for equivariant induction} proves the topology and cutoffs
in NCF.1--NCF.2, scalar Radon representation and Fubini in NCF.3, Haar
existence and translation continuity in NCF.4, Banach and strict
operator integration in NCF.5, and the locally supported subgroup
average in NCF.6.

\subsection{1. Actions and continuity}\label{actions-and-continuity}

An action \(\alpha\) on a C*-algebra \(A\) is point-norm continuous:
\(g\mapsto\alpha_g(a)\) is norm continuous for each \(a\). On a Hilbert
\(B\)-module \(E\), an action \(U\) consists of even complex-linear maps
satisfying \[
\begin{gathered}
U_g(\xi b)=U_g\xi\,\beta_g(b),\\
\langle U_g\xi,U_g\eta\rangle
 =\beta_g(\langle\xi,\eta\rangle),\\
U_gU_h=U_{gh},\qquad U_e=1.
\end{gathered}
\tag{1.1}
\] We require \(g\mapsto U_g\xi\) to be norm continuous for each vector.
The maps \(U_g\) are isometries, although they are generally not
\(B\)-linear. Conjugation defines an action on adjointable \(B\)-linear
operators: \[
g\cdot T=U_gTU_g^{-1}.
\tag{1.2}
\] An operator is \textbf{G-continuous} when its orbit under (1.2) is
norm continuous.

\textbf{Lemma 1.1 (compact operators and strict integration).} Compact
module operators are G-continuous. Every adjointable operator has a
strictly continuous orbit. If \(G\) is compact, a bounded strictly
continuous orbit can be integrated to an adjointable operator; the
average of an operator is invariant.

\textbf{Proof.} For a rank-one operator, \[
g\cdot\theta_{\xi,\eta}
 =\theta_{U_g\xi,U_g\eta}.
\tag{1.3}
\] The estimate \(\|\theta_{\xi,\eta}\|\leq\|\xi\|\|\eta\|\) proves norm
continuity. Finite sums and norm approximation give it for every compact
operator.

For \(T\in\mathcal L(E)\), the vector \(U_gT U_g^{-1}\xi\) depends
continuously on \(g\): use continuity of \(U_g^{-1}\xi\), boundedness of
\(T\), and continuity of \(U_g\) on each fixed vector. The same holds
for \(T^*\). On a bounded family, this strong continuity for the
operator and its adjoint is precisely strict continuity in
\(M(\mathcal K(E))\).

For compact \(G\), define the average on each vector by its Bochner
integral. It is bounded and \(B\)-linear, and integrating the adjoints
gives its adjoint. A change of Haar variable proves invariance.
Positivity and the contraction bound pass to this integral. If an
integrand is norm continuous with values in compact operators, its
integral is compact, since Riemann sums converge in norm. \(\square\)

A Kasparov G-cycle is a triple \((E,\phi,F)\), with an equivariant
graded homomorphism \(\phi:A\to\mathcal L(E)\) and an odd G-continuous
\(F\), such that \[
\begin{gathered}
{}[F,\phi(a)]_{\mathrm{gr}}\in\mathcal K(E),\\
(F^2-1)\phi(a)\in\mathcal K(E),\\
(F-F^*)\phi(a)\in\mathcal K(E),\\
(g\cdot F-F)\phi(a)\in\mathcal K(E)\\
(a\in A,\ g\in G).
\end{gathered}
\tag{1.4}
\] The fourth defect is norm continuous in \(g\), because \(F\) is
G-continuous. A cycle is degenerate if every defect in (1.4) is zero. A
homotopy is a cycle with coefficient \(C([0,1],B)\), with the given
action on \(B\) and trivial action on the interval. Its homotopy group
is \(KK^G(A,B)\). We set \[
KK_G^1(A,B)=KK^G(A,B\widehat\otimes C_1),
\tag{1.5}
\] where \(G\) acts trivially on the auxiliary Clifford algebra. We also
write \(KK_G^0=KK^G\), \(K_i^G(B)=KK_G^i(\mathbb C,B)\), and
\(K_G^i(A)=KK_G^i(A,\mathbb C)\).

The inverse-cycle rotation, direct-sum rotation and degenerate-cycle
contraction from \emph{Kasparov modules and the groups KK(A,B)} are
equivariant: their scalar rotation matrices commute with \(G\), and the
degenerate-cycle contraction is the C\_0-path proof of Proposition 2.3
of Lesson 06, with the diagonal group action. The action is strongly
continuous on these sections, by approximation by finite sums of scalar
functions times module vectors. These constructions prove that the
homotopy classes form an abelian group. Equivariant pullback and
coefficient tensoring give the two functorialities, with the action
\(U_g(\xi\widehat\otimes b)=U_g\xi\widehat\otimes\beta_g(b)\). The
rank-one and creation formulas of that lesson prove preservation of
compact defects; (1.4) tensors in the same way.

Self-adjoint normalization and clipping work for any locally compact
group: polynomial functional-calculus approximation preserves the
localized fourth defect, and the operator orbit remains norm continuous.
Compactness is needed for the invariant averaging in the next
proposition.

\textbf{Proposition 1.2 (compact-group normalization).} If \(G\) is
compact, every cycle and homotopy can be replaced, through an
equivariant locally compact perturbation, by one whose operator is
invariant, self-adjoint and contractive.

\textbf{Proof.} First use the self-adjoint normalization and clipping
from Proposition 3.1 and Theorem 3.2 of \emph{Kasparov modules and the
groups KK(A,B)}. Equivariance modulo compacts passes to these
operations. For clipping, in the quotient after source localization the
images of \(g\cdot F\) and \(F\) agree, so polynomial approximation to
the clipping function gives the same equality. Norm continuity of the
orbit passes to continuous functional calculus.

Average the normalized operator: \[
\overline F=\int_G g\cdot F\,dg.
\tag{1.6}
\] For each \(a\), \((\overline F-F)\phi(a)\) is the integral of the
fourth defect in (1.4), hence compact. The integrand is norm continuous
because \(F\) is G-continuous; thus this is a compact-operator norm
integral, not merely a strict-limit assertion. Taking adjoints gives
compactness on the other side. Put \(D=\overline F-F\) and \(F_t=F+tD\).
Its square satisfies \[
(F_t^2-1)\phi(a)=(F^2-1)\phi(a)
 +t(FD+DF)\phi(a)+t^2D^2\phi(a).
\] The \(FD\) and \(D^2\) terms are compact because \(D\phi(a)\) is
compact. For homogeneous \(a\), moving \(F\) through \(\phi(a)\) in the
\(DF\) term leaves \((-1)^{|a|}D\phi(a)F\) and a compact commutator, so
that term is compact too. The commutator of \(D\) with \(\phi(a)\) is
compact because both products are compact; the adjoint defect is zero.
The fourth defect is exactly \((1-t)(g\cdot F-F)\phi(a)\), and the orbit
of \(F_t\) is norm continuous. Hence this is an equivariant operator
homotopy. Haar translation makes the endpoint invariant; averaging
preserves its self-adjointness and contraction bound. Applying these
same equations over \(C([0,1],B)\) proves the homotopy assertion.
\(\square\)

\subsubsection{A two-element action as a working
model}\label{a-two-element-action-as-a-working-model}

Before constructing a product, consider an action of \(G=\{e,s\}\),
\(s^2=e\), on \(B\), with \(\alpha=\beta_s\) and \(\alpha^2=1\). On
\(B\oplus B\) use inner product \(\frac12(x_e^*y_e+x_s^*y_s)\),
reflecting normalized Haar measure. The group acts by \[
U_s(x_e,x_s)=(\alpha(x_s),\alpha(x_e)).
\] An operator matrix with entries in \(B\) is invariant precisely when
it has the form \[
\begin{pmatrix}A&C\\\alpha(C)&\alpha(A)\end{pmatrix}.
\] Indeed conjugating a matrix \(\begin{pmatrix}A&C\\D&E\end{pmatrix}\)
by \(U_s\) gives
\(\begin{pmatrix}\alpha(E)&\alpha(D)\\\alpha(C)&\alpha(A)\end{pmatrix}\).
Matrices over \(B\) are exactly the compact operators on this finite
free module, also for nonunital \(B\): products \(xy^*\) span a dense
subspace of each entry. The scalar factor in the inner product is
removed by the module unitary \((x_e,x_s)\mapsto2^{-1/2}(x_e,x_s)\).

Write a normalized-Haar coefficient function as \(f(e)=a,f(s)=b\). Its
integrated regular operator is \[
T_f=\frac12
\begin{pmatrix}a&b\\\alpha(b)&\alpha(a)\end{pmatrix}.
\tag{1.7}
\] Thus the counting-measure coefficients are \(a/2,b/2\). For
\(g(e)=c,g(s)=d\), direct multiplication gives \[
\begin{aligned}
(f*g)(e)&=\tfrac12(ac+b\alpha(d)),\\
(f*g)(s)&=\tfrac12(ad+b\alpha(c)),\\
f^*(e)&=a^*,\\
f^*(s)&=\alpha(b^*).
\end{aligned}
\tag{1.8}
\] These formulas give \(T_{f*g}=T_fT_g\) and \(T_{f^*}=T_f^*\).
Conversely every displayed invariant compact matrix is \(T_f\), with
coefficients twice its first row. This is already the concrete algebraic
content of the kernel identification in Proposition 4.1. That
proposition will prove that the same identification respects the full
crossed-product norm for every compact group.

For \(B=\mathbb C\) with trivial action, the constant orthonormal change
of coordinates to \((1,1)\) and \((1,-1)\) diagonalizes (1.7). Its two
entries are \((a+b)/2\) and \((a-b)/2\). The normalized coefficient
functions \((1,1)\) and \((1,-1)\) are therefore the two orthogonal
central projections; the identity function in this convolution algebra
is \((2,0)\). The two scalar blocks have ordinary K-groups
\(K_0=\mathbb Z^2\), \(K_1=0\), by the matrix computations in the
operator K-theory prerequisites. Theorem 4.3 will identify these blocks
with equivariant Fredholm classes. Proposition 6.1 will identify their
labels with the trivial and sign representations, with the tensor rule
\(\mathrm{sign}\otimes\mathrm{sign}=\mathrm{trivial}\).

Two different tasks remain visible in this small model. Invariance
constrains an operator's entries, while a Kasparov product must also
satisfy its connection and positivity conditions. Averaging an arbitrary
operator matrix settles only the first task. The invariant partition
construction below handles both tasks together. Later, stabilization
lets the finite matrix picture describe cycles on arbitrary countably
generated modules, and the kernel theorem replaces this finite
calculation by one over a compact group.

\subsection{2. Invariant partitions and the product for compact
groups}\label{invariant-partitions-and-the-product-for-compact-groups}

\textbf{Theorem 2.1 (the invariant technical theorem).} Let a compact
group act on a graded \(\sigma\)-unital algebra \(J\). Suppose
\(A_1,A_2\subset M(J)\) are invariant graded \(\sigma\)-unital
subalgebras, their elements are G-continuous, and \(\Delta\) is an
invariant separable graded subspace of G-continuous multipliers. Assume
\[
A_1A_2\subset J,\qquad
[\Delta,A_1]_{\mathrm{gr}}\subset A_1.
\tag{2.1}
\] There are invariant even positive contractions \(M,N\) such that \[
\begin{gathered}
M+N=1,\\
MA_1\subset J,\qquad NA_2\subset J,\\
[M,\Delta]\subset J.
\end{gathered}
\tag{2.2}
\] Their square roots have the annihilation and commutator properties of
Corollary 3.2 of \emph{Kasparov's technical theorem}.

\textbf{Proof.} Theorem 3.1 of that lesson first supplies \(M_0,N_0\),
without invariance. The action on \(M(J)\) is strictly continuous on
each fixed multiplier, by the same argument as Lemma 1.1 with module
\(J\). Take strict Haar averages \(M,N\).

For \(a\in A_1\), the function \[
\begin{gathered}
g\longmapsto(g\cdot M_0)a\\
=g\cdot\bigl(M_0(g^{-1}\cdot a)\bigr).
\end{gathered}
\tag{2.3}
\] is norm continuous with values in \(J\). Membership uses invariance
of \(A_1\); continuity uses norm continuity of \(g^{-1}\cdot a\), fixed
bounded multiplication by \(M_0\), and norm continuity of the action on
\(J\). Therefore \(Ma\in J\). The argument for \(NA_2\) is identical.
For \(d\in\Delta\), \[
[g\cdot M_0,d]
 =g\cdot[M_0,g^{-1}\cdot d]
\tag{2.4}
\] is another norm-continuous \(J\)-valued function. Its integral proves
the commutator conclusion. Positivity, evenness and \(M+N=1\) survive
averaging. Invariance follows by Haar translation. Polynomial
approximation in the quotient proves the assertions about roots, exactly
as in the cited corollary. \(\square\)

\textbf{Lemma 2.2 (invariant connections).} For compact \(G\), let
\(E_1,E_2\) be equivariant modules and let \(F_2\) be invariant and
self-adjoint. On \(E_1\widehat\otimes_D E_2\), there is an invariant
self-adjoint odd \(F_2\)-connection.

\textbf{Proof.} Choose a non-equivariant self-adjoint connection \(G_0\)
by the Grassmann connection construction in Section 1 of
\emph{Connections and the existence of the Kasparov product}. The
tensor-module action is diagonal. For homogeneous \(\xi\in E_1\),
conjugation carries its creation operator to the creation operator of
\(U_g\xi\). Thus the creation error of \(g\cdot G_0\) for \(\xi\) is the
conjugate of the creation error of \(G_0\) for \(U_g^{-1}\xi\). This
error depends continuously on \(g\) in the norm of compact operators
between the two modules: creation operators depend norm continuously on
their vectors, and conjugation on compact operators is norm continuous
by (1.3). The adjoint creation error has the same property. Strictly
average \(G_0\). The two integrated errors are compact, so its invariant
average is a connection. Strict averaging preserves its parity and
self-adjointness. \(\square\)

\textbf{Theorem 2.3 (compact-group composition).} For compact \(G\),
composition and external product give bilinear natural pairings \[
\begin{gathered}
KK_G^i(A,D)\times KK_G^j(D,B)\\
\longrightarrow KK_G^{i+j}(A,B),\\[3pt]
KK_G^i(A_1,B_1)\times KK_G^j(A_2,B_2)\\
\longrightarrow KK_G^{i+j}(A_1\widehat\otimes A_2,\\
 B_1\widehat\otimes B_2).
\end{gathered}
\tag{2.5}
\] Composition is associative, has the homomorphism cycles as
identities, and respects equivariant homomorphisms. External product has
the signed flip symmetry. Equivariant module homotopy and stable
equivariant operator homotopy give the same groups.

\textbf{Proof.} Normalize both input operators invariantly by
Proposition 1.2. Put \(S=F_1\widehat\otimes1\) and choose the invariant
connection \(G\) from Lemma 2.2. In the notation of the non-equivariant
existence proof, put \[
\begin{gathered}
I_0=\mathcal K(E_1)\widehat\otimes1,\\
J=\mathcal K(E_1\widehat\otimes_D E_2),\\
A_1=C^*(I_0,J),\\
\begin{aligned}
A_2=C^*\bigl(&J,G^2-1,[S,G]_{\mathrm{gr}},\\
 &[G,\phi(A)]_{\mathrm{gr}}\bigr),
\end{aligned}\\
\Delta=\overline{\operatorname{span}}\\
\{S,G,\phi(A),S^*,G^*,\phi(A)^*\}.
\end{gathered}
\tag{2.6}
\] These are invariant: \(S,G\) are invariant, \(I_0,J\) are invariant
ideals in the indicated represented algebra, and equivariance of
\(\phi\) preserves the last family. All their elements are G-continuous,
by norm approximation from these generators. Separability of \(A\) and
countable generation give the countable tests and strictly positive
elements constructed in the non-equivariant proof. That proof's
zero-connection calculation proves \(A_1A_2\subset J\), and its
connection calculus proves that \(\Delta\) derives \(A_1\).

Apply Theorem 2.1. The invariant operator \[
F=M^{1/4}SM^{1/4}+N^{1/4}GN^{1/4}
\tag{2.7}
\] is self-adjoint and odd. Equations (3.5)--(3.7) of the
non-equivariant product lesson prove its localized square defect,
representation commutator, connection condition and positivity. Each of
those equations involves precisely the data (2.6) and the partition
(2.2); therefore it applies here without any changed analytic
assumption. Invariance supplies the additional fourth condition (1.4)
with zero defect.

An initially noninvariant equivariant product operator can also be
averaged. Its connection errors average to compact errors by Lemma 2.2,
and its localized anticommutator with invariant \(S\) averages to a
positive quotient element. Proposition 1.2 gives its locally compact
homotopy. Thus uniqueness can be checked among invariant operators.

For uniqueness, use the same invariant \(A_1\), and replace \(A_2\) by
\[
C^*(J,[S,F]_{\mathrm{gr}},[S,F']_{\mathrm{gr}},F-F').
\tag{2.8}
\] All generators are invariant zero-connections. Include the invariant
\(S,F,F'\) and equivariant \(\phi(A)\) in \(\Delta\). Theorem 2.1 then
gives an invariant comparison operator. Equation (4.3) and the
positive-anticommutator homotopy in that lesson join both product
operators to it through invariant operators. The argument also works
over \(C([0,1],B)\), so product descends to homotopy classes in both
variables. Block sums prove bilinearity.

Essential replacement, when a homomorphism is an input, uses the
interval ideal and module of Proposition 5.1 of that lesson. They are
invariant under the pointwise action. Choose their connection by Lemma
2.2. The compact creation calculations in that proposition then give an
equivariant interval cycle. Its endpoint tensor unitaries are
equivariant. Thus Theorem 5.2 there proves both homomorphism identities
here.

For associativity, use Lemmas 3.1--3.2 and the explicit separation proof
of Theorem 4.1 in \emph{Homotopy, associativity, the index pairing and
KK-equivalence}. Its bracket connections and their negative parts are
equivariant, because all selected input and intermediate operators are
invariant and functional calculus commutes with conjugation. Its
separation algebras, generated by those operators, their compact defects
and \(\phi(A)\), are invariant. Replace the single technical-theorem
step in that proof by Theorem 2.1. The resulting simultaneous product
operator is invariant and is a product for each parenthesization by the
exact positivity estimates there. The invariant uniqueness just proved
identifies the two classes.

The evaluation bridge in Lemma 1.1 of the homotopy lesson has trivial
action on its scalar Hilbert spaces. Its exterior extension with \(B\),
and all its scalar rotation paths, are equivariant. Using the just
established equivariant products and uniqueness, the proof of Theorem
2.3 there converts every module homotopy into a stable equivariant
operator homotopy. This proves the last assertion, without assuming
equivariant stabilization.

Finally, external product is formed by exterior extension and
composition with diagonal actions. Its signed flip unitaries are
equivariant. The two-order comparison of Theorem 5.3 of the homotopy
lesson now uses the established invariant associativity and proves
signed symmetry. The degree-one construction uses the same trivial
auxiliary Clifford factors and the same ordered odd--odd Morita map as
in that theorem. This also proves every parity of (2.5). \(\square\)

\subsubsection{Almost invariant
partitions}\label{almost-invariant-partitions}

Exact averaging is replaced by uniform approximate invariance on compact
sets for a noncompact group.

\textbf{Lemma 2.4 (convex approximate invariance).} Let a locally
compact group act continuously on a \(\sigma\)-unital algebra \(D\).
Given a compact \(K\subset G\), finite or norm-compact
derivating-multiplier commutator tests and finite approximate-identity
tests, finite convex combinations of any sufficiently late tail of a
nested functional-calculus approximate identity can satisfy all the
tests and \[
\sup_{g\in K}\|g\cdot v-v\|<\varepsilon.
\tag{2.9}
\] Successive choices can retain \(v_nv_{n-1}=v_{n-1}\), and can be even
for graded data.

\textbf{Proof.} Take the nested \(e_j=f_j(h)\) of Section 1 of
\emph{Kasparov's technical theorem}. In \(C(K,D)\), both functions
\(g\mapsto e_j\) and \(g\mapsto g\cdot e_j\) are positive contractive
approximate identities. For the second assertion and \(z\in C(K,D)\), \[
\begin{gathered}
\|(g\cdot e_j)z(g)-z(g)\|\\
=\|e_j(g^{-1}\cdot z(g))-g^{-1}\cdot z(g)\|.
\end{gathered}
\tag{2.10}
\] The right-hand test values form a norm-compact subset of \(D\), so
convergence is uniform; the other-side test is identical. The
bounded-functional construction in
\href{KT-KK-08.html\#1-approximate-identities-with-commutator-control}{Lesson
08, Theorem 1.1} writes every bounded functional on \(C(K,D)\) as a
vector coefficient of a representation. Both represented approximate
identities converge strongly to the same projection onto the closure of
the represented algebra acting on that Hilbert space, by the
strong-support argument proved there. Their represented difference tends
strongly to zero; testing each such vector coefficient proves weak
convergence to zero in the full Banach dual of \(C(K,D)\).

Combine that difference with the finite commutator tuple from Theorem
1.1 of the technical-theorem lesson. The tuple has weak limit zero in
the finite direct sum. Hahn--Banach separation gives a single finite
convex combination making every norm small, including the sup norm in
(2.9). Compact commutator sets reduce to finite nets. A sufficiently
late tail already satisfies the finite approximate-identity tests, and
convex combinations retain their bounds. Choosing the next indices
beyond the previous finite support gives \(v_nv_{n-1}=v_{n-1}\) by the
nested product equation. Evenness is preserved throughout. \(\square\)

\textbf{Theorem 2.5 (the locally compact technical theorem).} Under the
invariant separation hypotheses of Theorem 2.1, but for any second
countable locally compact \(G\), there are even G-continuous positive
contractions satisfying (2.2) and \[
g\cdot M-M\in J\quad(g\in G),
\tag{2.11}
\] with this difference norm continuous in \(g\).

\textbf{Proof.} For separable tested algebras retain every test in the
strict-series proof of Theorem 2.1 of \emph{Kasparov's technical
theorem}. Its approximate identities are \(u_n\in A_1\), \(v_n\in J\),
its square-root differences are \(b_n=(v_n-v_{n-1})^{1/2}\), and its sum
is \[
\begin{gathered}
N=\sum_{n\geq1}b_nu_nb_n,\\
M=1-N.
\end{gathered}
\tag{2.12}
\] Equations (2.3)--(2.12) of that proof establish separation and
commutation modulo \(J\). Lemma 2.4 allows the additional group tests
simultaneously with those original tests.

Choose increasing compact \(K_n\) whose interiors cover \(G\), and
\(\varepsilon_n=2^{-n}\). Require \[
\begin{gathered}
\sup_{g\in K_n}\|g\cdot u_n-u_n\|<\varepsilon_n,\\
\sup_{g\in K_n}\|g\cdot b_n-b_n\|<\varepsilon_n.
\end{gathered}
\tag{2.13}
\] For the second requirement, uniform polynomial approximation to the
square root supplies a tolerance \(\eta_n\). At step \(n\) require the
group difference of \(v_n\), over \(K_{n+1}\), to be smaller than half
both \(\eta_n\) and \(\eta_{n+1}\). Then the already chosen \(v_{n-1}\)
and the new \(v_n\) give the desired bound for the square root of their
difference over \(K_n\). The first difference, with \(v_0=0\), has its
own bound. This is the prospective square-root selection of the
preceding proof, now with both commutator and group tests.

For \(g\in K_n\), expansion into three differences gives \[
\begin{gathered}
\|g\cdot(b_nu_nb_n)-b_nu_nb_n\|\\
\leq3\varepsilon_n.
\end{gathered}
\tag{2.14}
\] Every difference is in \(J\). On fixed \(K_m\) the tail with
\(n\geq m\) converges uniformly in norm; the finitely many initial terms
have continuous orbits in \(J\). The difference series is therefore norm
continuous and \(J\)-valued. Strict continuity of the action identifies
its sum with \(g\cdot N-N\): apply the action to bounded strict partial
sums and test on elements of \(J\). This proves (2.11) and G-continuity.

For general \(\sigma\)-unital separation algebras, use the separable
reduction at the end of the preceding technical-theorem proof, also
closing its subalgebras under translates. Begin with strictly positive
\(h_0,h_1,h_2\), include their orbits, and close under the required
products and derivations. A countable dense subset of \(G\) suffices
because the tested orbits are norm continuous. The resulting
\(J_0,B_1,B_2\) are separable and invariant. The proved embedding
\(M(J_0)\subset M(J)\) is equivariant by strict extension from invariant
\(J_0\). The separation conclusions extend using the original \(h_i\);
the norm-continuous difference series remains in \(J_0\subset J\).
Graded invariant subalgebras and even approximate identities retain
parity. \(\square\)

Invariance of the separation data matters: an almost invariant partition
cannot separate two supports that the group exchanges modulo \(J\). The
product data below satisfy the invariant hypothesis.

\textbf{Theorem 2.6 (composition for locally compact groups).} The
product, naturality, associativity, signed symmetry and
homotopy-comparison conclusions of Theorem 2.3 hold for any second
countable locally compact \(G\).

\textbf{Proof.} First construct a G-continuous connection. Choose a
self-adjoint non-equivariant connection \(G_0\) and a nonnegative
\(f\in C_c(G)\) of integral one. Its strict convolution is \[
G_f=\int_G f(g)(g\cdot G_0)\,dg.
\tag{2.15}
\]
\href{supporting/noncompact-foundations/noncompact-foundations.html\#ncf-004}{NCF.4--NCF.5}
prove Haar existence, the strict integral and translation continuity of
\(f\) in \(L^1(G)\). They give norm continuity of the orbit, with bound
\(\|G_0\|\|\lambda_hf-f\|_1\). A translate of \(G_0\) is a connection
for \(g\cdot F_2\). The latter differs two-sided locally compactly from
\(F_2\), so the creation compactness test of Lemma 2.1 of the product
lesson makes it also an \(F_2\)-connection. For a fixed creation vector
its errors are norm continuous in \(g\): transformed vectors are norm
continuous, compact errors transform norm continuously, and \(F_2\) has
a norm-continuous orbit. Integrating the errors over compact support
proves that \(G_f\) is a self-adjoint odd connection. Subtracting its
translated connection identities shows that \(g\cdot G_f-G_f\) is a
zero-connection.

Normalize the inputs and set \(G=G_f\). Keep invariant \(A_1=I_0+J\)
from (2.6). Generate \(A_2\) by every translate of the defect generators
there and of \(g\cdot G-G\). These are zero-connections, hence
\(A_1A_2\subset J\). Include the orbits of \(S,G,\phi(A)\) and their
adjoints in \(\Delta\). The connection calculus makes them derivations
of invariant \(A_1\). Norm-continuous orbits generate separable tested
data; the positive generator-sum argument in the non-equivariant proof
gives \(\sigma\)-unitality. Apply Theorem 2.5.

The non-equivariant equations (3.5)--(3.7) prove all product conditions
for \(F_0=M^{1/2}S+N^{1/2}G\). Its additional defect is \[
\begin{gathered}
g\cdot F_0-F_0\\
 =(g\cdot M^{1/2}-M^{1/2})(g\cdot S)\\
 \quad+M^{1/2}(g\cdot S-S)\\
 \quad+(g\cdot N^{1/2}-N^{1/2})(g\cdot G)\\
 \quad+N^{1/2}(g\cdot G-G).
\end{gathered}
\tag{2.16}
\] The first and third terms are globally compact by (2.11) and quotient
functional calculus. Multiplying the second by \(\phi(a)\) gives a
compact operator, since \((g\cdot S-S)\phi(a)\in I_0\) and \(M^{1/2}\)
annihilates \(I_0\) modulo \(J\). The last term is globally compact by
the extra \(A_2\) generators. All factors are G-continuous. Thus (1.4)
holds, and the self-adjoint sandwich (2.7), differing globally compactly
from \(F_0\), has the same equivariant class.

For uniqueness start with (2.8), adjoin \(g\cdot F-F\),
\(g\cdot F'-F'\), and close under translates. These differences are
zero-connections because the translated second-input operators differ
locally compactly from \(F_2\). Include every tested operator orbit in
the derivating space. The comparison operator satisfies (1.4) by (2.16).
Its positive-anticommutator homotopies use continuous functional
calculus in G-continuous operators; polynomial approximation preserves
their localized equivariance defects. Hence uniqueness and the
interval-coefficient proof of bilinearity give actual equivariant
homotopies.

The associativity and evaluation-bridge arguments used in Theorem 2.3
have the same additional tests at each separation step. Their first
separation algebra is invariant: it is generated by the relevant tensor
images of compact operators and \(J\). Close the defect and derivating
sets under translates and include the equivariance differences of
connection operators being combined. The added defects annihilate the
invariant first algebra modulo \(J\), by zero-connection calculus or by
translating the original annihilation equation. The added derivating
tests still derive it. Theorem 2.5 supplies the partition; expansion
(2.16), with the corresponding localized first-module defect, proves
equivariance of the simultaneous product or homotopy operator. The
original connection and positivity equations are unchanged. The exact
preceding proofs therefore give associativity, essential replacement,
homomorphism identities and comparison of homotopies. The external
signed flips commute with the diagonal group action. This proves all
asserted properties without averaging over a noncompact group.
\(\square\)

\subsection{3. Equivariant
stabilization}\label{equivariant-stabilization}

Write \[
\begin{gathered}
\mathcal H_B^G=L^2(G)\otimes\ell^2\otimes B,\\
(U_g\zeta)(s)=\beta_g(\zeta(g^{-1}s)).
\end{gathered}
\tag{3.1}
\] For grading, give \(\ell^2\) infinitely many even and odd coordinates
and use the tensor grading. This module has a strongly continuous action
by approximation by continuous functions with finitely many coordinates.

\textbf{Theorem 3.1 (compact equivariant stabilization).} Assume \(G\)
is compact. Every countably generated graded equivariant Hilbert
\(B\)-module \(E\) admits an even equivariant adjointable isometry into
\(\mathcal H_B^G\), and \[
E\oplus\mathcal H_B^G\cong\mathcal H_B^G
\tag{3.2}
\] by an even equivariant module unitary.

\textbf{Proof.} Choose homogeneous generating vectors \(\xi_j\) of norm
at most one and positive weights \(a_j\) with \(\sum a_j^2<\infty\).
Include enough vectors that \(k_0=\sum a_j^2\theta_{\xi_j,\xi_j}\) is
strictly positive in \(\mathcal K(E)\). Here is a direct check. It
dominates \(a_j^2\theta_{\xi_j,\xi_j}\). If \(e_n\) is its spectral
cutoff approximate identity, then
\(\|(1-e_n)\xi_j\|^2\leq a_j^{-2}\|(1-e_n)k_0(1-e_n)\|\to0\). The
generating property extends this to every vector. Applying it to both
vectors in each rank-one operator proves that \(e_n\) is an approximate
identity of \(\mathcal K(E)\).

Assign the \(j\)-th regular coordinate the parity of \(\xi_j\). Define
\[
\begin{gathered}
T\zeta=\sum_j a_j\int_G U_s\xi_j\,\zeta_j(s)\,ds,\\
(T^*\eta)_j(s)=a_j\langle U_s\xi_j,\eta\rangle.
\end{gathered}
\tag{3.3}
\] On finite continuous sections the formulas are adjoint to each other.
The Hilbert-module Cauchy--Schwarz inequality for the integral gives
\(\|T_j\|\leq\|\xi_j\|\); summability of \(a_j^2\) gives a bounded row
operator and its adjoint column. The formulas extend by density. They
are \(B\)-linear; changing \(s=gr\) in the first proves \(TU_g=U_gT\).
In particular \[
k=TT^*=\int_G g\cdot k_0\,dg
\tag{3.4}
\] is invariant and compact.

It is strictly positive in \(\mathcal K(E)\). Indeed every nonzero
positive functional \(\omega\) on that algebra has \(\omega(k_0)>0\).
The continuous nonnegative function \(g\mapsto\omega(g\cdot k_0)\) is
positive at the identity, so its integral is positive. The
characterization of strict positivity by positive functionals proves the
claim. Equivalently, its spectral cutoff approximate identity tends
strictly to one, so \(kE\) is dense in \(E\).

Choose a nonnegative locally finite continuous partition \(\chi_n\) of
\((0,\|k\|]\), with each support bounded away from zero. Set
\(f_n(t)=(\chi_n(t)/t)^{1/2}\), extended by zero at zero. Then \[
\begin{gathered}
V\eta=(T^*f_n(k)\eta)_n,\\
\sum_n f_n(k)k f_n(k)=1\quad\text{strictly}.
\end{gathered}
\tag{3.5}
\] Consequently \(V\) is an isometry into a countable direct sum of
regular modules. It is adjointable: the partial row operators
\(\sum_n f_n(k)T\zeta_n\) have norm at most one, and convergence on
finitely supported vectors extends them to the adjoint. Functional
calculus of invariant even \(k\) makes \(V\) equivariant and even. The
resulting countable sum is another copy of \(\mathcal H_B^G\).

Write this copy as \(VE\oplus E^\perp\). Its countable direct sum is
\(E^\infty\oplus(E^\perp)^\infty\). Absorbing one additional copy of
\(E\) into \(E^\infty\) by the coordinate shift gives (3.2). All
coordinate rearrangements are even and equivariant. \(\square\)

The averaged strictly positive operator is the reason this argument
gives an exactly equivariant unitary for compact groups. An integral
over all of a noncompact group would require a separate construction.

\textbf{Theorem 3.2 (G-continuous stabilization).} For any second
countable locally compact \(G\), there is an even G-continuous module
unitary \[
E\oplus\mathcal H_B^G\longrightarrow\mathcal H_B^G.
\tag{3.6}
\] It need not be equivariant.

\textbf{Proof.} Choose a unit vector \(f\in C_c(G)\subset L^2(G)\).
Separate \(f\otimes E^\infty\) from its orthogonal complement in
\(L^2(G,E^\infty)\). Send an additional copy of \(E\) to the zeroth
\(f\)-coordinate, shift the old \(f\)-coordinates by one, and fix the
orthogonal complement. This gives an even unitary \[
\begin{gathered}
W:E\oplus L^2(G,E^\infty)\\
\longrightarrow L^2(G,E^\infty).
\end{gathered}
\tag{3.7}
\] Its components are the isometry \(\xi\mapsto f\otimes\xi\), the
projection \(p_f\), and scalar coordinate shifts. Conjugation replaces
\(f\) by \(\lambda_gf\) and \(p_f\) by \(p_{\lambda_gf}\); the action on
\(E\) cancels. The bounds \(\|\lambda_gf-f\|_2\to0\) and
\(\|p_{\lambda_gf}-p_f\|\leq2\|\lambda_gf-f\|_2\) prove G-continuity.

For an even module unitary \(w:E_1\to E_2\), regularize it by \[
(\widetilde w\xi)(t)=
 U_t^{(2)}w(U_t^{(1)})^{-1}\xi(t).
\tag{3.8}
\] The semilinear coefficient actions cancel, giving \(B\)-linearity.
The field and its adjoint act continuously on compact continuous
sections by strong continuity of both actions, and extend to inverse
L²-module unitaries. Substitution of \(g^{-1}t\) proves exact
equivariance.

Ordinary graded stabilization gives an even unitary
\(E^\infty\oplus\widehat H_B\cong\widehat H_B\). Regularize it,
obtaining an equivariant unitary
\(L^2(G,E^\infty)\oplus\mathcal H_B^G\cong\mathcal H_B^G\). Use its
inverse to open the regular summand in (3.6), apply \(W\oplus1\), and
close with the same equivariant unitary. The composite is G-continuous.
All L² modules are countably generated by second countability.
\(\square\)

\subsection{4. Invariant compact operators and
Green--Julg}\label{invariant-compact-operators-and-greenjulg}

For the next calculation use the unitarily equivalent regular-module
action \[
\begin{gathered}
\mathcal R_B=L^2(G)\otimes B,\\
(U_g\xi)(t)=\beta_g(\xi(tg)).
\end{gathered}
\tag{4.1}
\] Inversion \(t\mapsto t^{-1}\) intertwines (3.1) and (4.1), because
compact groups are unimodular. The regular covariant pair on this module
is \[
\begin{gathered}
(\pi(b)\xi)(t)=\beta_{t^{-1}}(b)\xi(t),\\
(\lambda_s\xi)(t)=\xi(s^{-1}t).
\end{gathered}
\tag{4.2}
\] Both operators in (4.2) commute with (4.1).

\textbf{Proposition 4.1 (the invariant kernel algebra).} For compact
\(G\), integration of (4.2) is an isomorphism \[
B\rtimes G\ \cong\ \mathcal K(\mathcal R_B)^G.
\tag{4.3}
\] It preserves grading. With an additional trivial infinite
multiplicity space it gives \[
(B\rtimes G)\widehat\otimes\mathcal K
 \cong\mathcal K(\mathcal H_B^G)^G.
\tag{4.4}
\]

\textbf{Proof.} An element \(f\in C(G,B)\) acts as the integral kernel
\[
K_f(t,r)=\beta_{t^{-1}}\bigl(f(tr^{-1})\bigr).
\tag{4.5}
\] Indeed substitute \(r=s^{-1}t\) in its integrated formula. This
kernel is norm continuous on compact \(G\times G\). Such a \(B\)-valued
kernel defines a compact module operator: approximate it uniformly by
finite sums \(a(t)b(r)c\), then approximate each coefficient \(c\in B\)
by sums of products \(uv^*\). The resulting kernels are rank-one
operators on \(L^2(G)\otimes B\). The operator-norm error is bounded by
the uniform kernel error, since Haar mass is one.

Conjugation in (4.1) sends a kernel \(K\) to \[
(g\cdot K)(t,r)=\beta_g(K(tg,rg)).
\tag{4.6}
\] Thus (4.5) is invariant. Conversely an invariant continuous kernel
obeys \(K(tg,rg)=\beta_{g^{-1}}K(t,r)\). Putting \(g=t^{-1}\) gives
(4.5) with \(f(s)=K(e,s^{-1})\). Continuous kernels are dense among
compact operators by approximation of both vectors of a rank-one
operator by continuous sections. Averaging that approximation in norm
makes it invariant, and preserves continuity of its kernel. Hence
invariant continuous kernels are dense in
\(\mathcal K(\mathcal R_B)^G\). This proves density of the range in
(4.3).

The regular representation is faithful on the full crossed product in
this compact case. To verify the full-to-reduced step, take any
nondegenerate covariant Hilbert-space pair \((\rho,V)\). The map
\(v\mapsto v\otimes1_G\) embeds it into
\((\rho\otimes1,V\otimes\lambda)\), since \(1_G\) is an invariant unit
vector. Fell absorption, with its explicit unitary from \emph{Reduced
crossed products and Fell's absorption principle}, identifies the latter
with the regular representation for \(\rho\). Regular-representation
norm independence in that lesson bounds it by the faithful regular norm.
Taking the supremum over all covariant pairs proves that full and
regular norms coincide. Tensoring the regular module with a faithful
Hilbert-space representation of \(B\) tests its norm faithfully;
therefore its integrated module representation is faithful as well. A
faithful homomorphism has closed range, so the already dense range is
all of the fixed compact algebra.

Finally compress compact operators to finitely many multiplicity
coordinates. Those compressions converge in norm and commute with the
group action. Their invariant entries are precisely the algebra in
(4.3). This proves (4.4), with the indicated tensor grading. \(\square\)

\textbf{Lemma 4.2 (fixed multipliers).} Put
\(J=\mathcal K(\mathcal H_B^G)^G\). Its action on \(\mathcal H_B^G\) is
nondegenerate, and strict extension identifies \[
M(J)=\mathcal L(\mathcal H_B^G)^G.
\tag{4.7}
\]

\textbf{Proof.} Average a strictly positive compact operator as in
(3.4), now on the regular module itself. The averaged operator lies in
\(J\) and remains strictly positive in the whole compact algebra. Its
cutoff approximate identity acts in norm on every vector. Thus \(J\)
acts nondegenerately, and the multiplier extension exists and is
faithful. Every extended multiplier is invariant, since its products
with \(J\) are invariant and those products determine its action on a
dense set of vectors. Conversely every invariant adjointable \(T\)
satisfies \(TJ,JT\subset J\), so is a multiplier of \(J\). The extension
agrees with \(T\) on the dense subspace \(J\mathcal H_B^G\), proving
surjectivity. \(\square\)

\textbf{Theorem 4.3 (Green--Julg for compact groups).} For compact
\(G\), there is a natural isomorphism \[
K_i^G(B)\cong K_i(B\rtimes G).
\tag{4.8}
\] For trivial grading on \(B\), the group on the right is ordinary
K-theory.

\textbf{Proof.} We give the cycle comparison, including its equivalence
relations. A scalar source has the invariant projection \(p=\phi(1)\).
Its commutator with \(F\) is compact. Removing the off-diagonal corners
of \(F\) is therefore a compact perturbation; the \(1-p\) corner has
zero representation and is degenerate. Thus take the scalar
representation to be unital. By Proposition 1.2 take \(F\) invariant.

Add a unital degenerate cycle on the balanced regular module, whose
operator is the odd flip between its even and odd multiplicity
coordinates. Theorem 3.1 identifies the enlarged module equivariantly
with the balanced regular module. The resulting cycle is an odd
invariant operator with adjoint and square defects in
\(J=\mathcal K(\mathcal H_B^G)^G\). Consequently it is exactly an odd
multiplier of \(J\), with those defects in \(J\), by Lemma 4.2.

On the other hand put \(D=B\rtimes G\). The ordinary scalar cycle
picture for \(KK(\mathbb C,D)\), after graded stabilization, consists of
odd operators on a balanced standard \(D\)-module with the same defects
in its compact algebra. Proposition 4.1 identifies that compact algebra
with \(J\), including grading; its isomorphism extends to multipliers.
It therefore identifies these two sets of normalized standard cycles.

It also identifies all their relations. Norm operator paths are the same
multiplier paths, compact perturbations are the same elements of \(J\),
and even unitary changes of standard-module coordinates are the same
even multiplier unitaries. A degenerate operator has exactly zero
defects on either side. Theorem 2.3 identifies equivariant module
homotopy with stable equivariant operator homotopy. The corresponding
non-equivariant identification is Theorem 2.3 of \emph{Homotopy,
associativity, the index pairing and KK-equivalence}. Hence the
standard-cycle identification descends to a bijection of the actual
homotopy groups, rather than merely their operator-path versions. Block
sums prove additivity.

For degree one apply the same argument with the trivial auxiliary
\(C_1\) action. The integrated kernel isomorphism then identifies
\((B\widehat\otimes C_1)\rtimes G\) with
\((B\rtimes G)\widehat\otimes C_1\). On continuous functions this is the
identity on the Clifford factor, and the full covariant universal
property gives its inverse. This proves (4.8) for either grading and
both degrees.

The construction is natural in equivariant coefficient maps. On a
rank-one kernel a coefficient map replaces each inner product by its
image; scalar extension of the regular module gives exactly that
replacement. Finite kernel approximation, then multiplier extension on
the generated support module, gives the same map for arbitrary cycles.
Thus the two normalized cycle constructions commute with coefficient
pushforward. \(\square\)

\subsection{5. The discrete dual}\label{the-discrete-dual}

The full crossed product is essential here. For a countable discrete
group its dense algebra consists of finite sums \(\sum_g a_g u_g\), and
\(a u_g\) belongs to the algebra even when \(A\) is nonunital. For a
nondiscrete group, an individual coefficient times a group multiplier
need not belong to the crossed product.

\textbf{Theorem 5.1 (the dual Green--Julg theorem).} Let \(G\) be
countable and discrete. There is a natural isomorphism \[
K_G^i(A)\cong K^i(A\rtimes G).
\tag{5.1}
\] Its forward map integrates the covariant representation and retains
the cycle operator.

\textbf{Proof.} First let the coefficient be \(\mathbb C\). For an
equivariant cycle \((H,\phi,F)\), integrate \((\phi,U)\) to
\(\Phi:A\rtimes G\to\mathcal L(H)\). For homogeneous \(a\), the graded
commutator expands as \[
\begin{gathered}
{}[F,\phi(a)U_g]_{\mathrm{gr}}\\
 = [F,\phi(a)]_{\mathrm{gr}}U_g\\
 \quad+(-1)^{|a|}\phi(a)[F,U_g].
\end{gathered}
\tag{5.2}
\] The second term is compact by the equivariance defect:
\([F,U_g]=(F-g\cdot F)U_g\), and the two-sided localized version of
(1.4) follows by taking adjoints after self-adjoint normalization. The
square and adjoint defects multiplied by \(\phi(a)U_g\) are compact by
the original cycle conditions. Finite sums and norm density extend these
statements to every element of \(A\rtimes G\). Thus \((H,\Phi,F)\) is a
non-equivariant cycle.

Conversely take a cycle \((H,\Phi,F)\) for \(D=A\rtimes G\). It can be
taken with nondegenerate representation. On a Hilbert space this also
follows directly by compressing to \(H_0=\overline{\Phi(D)H}\): the
orthogonal complementary representation is zero, and
\((1-P)F\Phi(d)=(1-P)[F,\Phi(d)]\) is compact. Removing these
off-diagonal terms is a two-sided locally compact perturbation, and the
complementary cycle is degenerate.

Extend the nondegenerate representation strictly to \(M(D)\). Define
\(\phi(a)=\Phi(a)\) and \(U_g=\widetilde\Phi(u_g)\). They are covariant.
The inclusion \(A\subset D\) is nondegenerate, so \(\phi\) is
nondegenerate. Every orbit is norm continuous because the group is
discrete. Both \([F,\Phi(au_g)]_{\mathrm{gr}}\) and
\([F,\phi(a)]_{\mathrm{gr}}\) are compact. Equation (5.2) therefore
proves \(\phi(a)[F,U_g]\) compact. Multiplying by \(U_g^{-1}\), and
taking adjoints for the other localization, gives the fourth defect in
(1.4). The remaining defects are obtained by restricting from \(D\).
Thus this is an equivariant cycle.

For completeness, the inverse respects homotopy even when a homotopy's
module is not a trivial Hilbert-space field. Given a \(D\)-\(C([0,1])\)
homotopy, use Proposition 5.1 of \emph{Connections and the existence of
the Kasparov product} to replace it by a cycle with nondegenerate
\(D\)-representation. The resulting essential module is
\(\overline{\Phi(D)E}\). Its representation extends to the canonical
\(u_g\), and (5.2), now in compact module operators over the interval,
supplies the equivariant cycle conditions. Evaluation has the essential
representation at each endpoint. The replacement operator at an endpoint
differs locally compactly from its compressed original operator: this is
the connection conclusion of the essential-replacement construction when
the original endpoint is essential. Such a difference gives an
equivariant straight homotopy, since the action is discrete and its
additional localized defect is compact. Therefore the inverse sends
homotopic cycles to the same equivariant class. This argument also
covers stable additions and degenerate cycles. Integration itself
preserves homotopies by exactly (5.2) on the interval module.

On nondegenerate covariant representations the two operations are
inverse by the integrated-form theorem. For a cycle with degenerate
source representation, its Hilbert-space support projection commutes
with \(U_g\), since the represented support is invariant. The preceding
compression is then equivariant. Thus the two operations are inverse on
all classes, not only essential representatives. Pulling a
representation back along an equivariant algebra map commutes with
integration, proving naturality.

For degree one use the ungraded odd Fredholm description of Proposition
3.1 of \emph{Pictures of KK}. The same covariant representation and the
same compact estimates apply. Equivalently retain the finite auxiliary
right \(C_1\) factor throughout; its action is trivial and its
multiplier extension causes no new continuity requirement. This proves
(5.1) in both degrees. \(\square\)

\subsection{6. Representation classes}\label{representation-classes}

For any group for which the equivariant product is constructed, define
\[
R^i(G)=KK_G^i(\mathbb C,\mathbb C).
\tag{6.1}
\] For compact groups Theorem 2.3 supplies its full ring structure.
Finite-dimensional unitary representations give even cycles \((V,1,0)\),
with \(V\) in even degree. Their square defect is compact because \(V\)
is finite dimensional.

\textbf{Proposition 6.1 (the compact representation ring).} For compact
\(G\), \[
R^0(G)=R(G),\qquad R^1(G)=0,
\tag{6.2}
\] where the right-hand ring is the Grothendieck ring of
finite-dimensional unitary representations. Addition is direct sum and
multiplication is tensor product. In particular it is commutative, and
it acts naturally on \(KK_G^i(A,B)\).

\textbf{Proof.} An invariant even Fredholm operator on a graded Hilbert
space has off-diagonal part \(T\) commuting with the group. Its polar
decomposition commutes with the group as well: form the support and
modulus by spectral calculus. Kernel and cokernel are finite-dimensional
invariant subspaces. The invertible complementary part is an equivariant
compact perturbation of its polar unitary and yields a degenerate cycle.
Thus every even class is \([
\ker T]-[\ker T^*]\). To check independence and injectivity, we describe
the invariant compact algebra and its K-groups rather than assume
invariance of these finite kernels under an arbitrary module homotopy.

Every strongly continuous unitary representation on a separable Hilbert
space is a Hilbert sum of finite-dimensional representations when \(G\)
is compact. Indeed average a positive compact operator with trivial
kernel, as in (3.4). The average still has trivial kernel: for a nonzero
vector its positive quadratic form is positive at the identity and hence
on a set of positive Haar measure. Its positive eigenspaces are finite
dimensional and invariant, and the compact spectral theorem gives their
dense orthogonal sum. A finite-dimensional unitary representation splits
into irreducibles by taking invariant orthogonal complements.

Every irreducible \(V\) occurs in the right regular representation: for
a nonzero linear functional \(\ell\) on \(V\), the map
\(v\mapsto(t\mapsto\ell(\pi(t)v))\) intertwines \(V\) with right
translation and is injective, since its kernel is invariant and the map
is nonzero. Schur's lemma follows directly by spectral decomposition of
a commuting self-adjoint matrix; a general commuting matrix has scalar
real and imaginary parts. Applying it to an averaged map between two
irreducibles shows that inequivalent irreducibles have zero intertwiner
space. Consequently, with infinite trivial multiplicity, \[
\mathcal K(\mathcal H_{\mathbb C}^G)^G
 \cong\bigoplus_{V\in\widehat G}^{c_0}
 \mathcal K(M_V)\otimes1_V,
\tag{6.3}
\] where each multiplicity space \(M_V\) is infinite dimensional. The
\(c_0\) condition follows from compactness: finitely many finite-rank
compressions approximate a compact operator, and its norm on all
remaining isotypical summands then tends to zero. Conversely a \(c_0\)
family of compact blocks tensored with the finite-dimensional identities
is compact. Second countability makes this a countable collection, as
each irreducible occurs in separable \(L^2(G)\).

The ordinary K-groups of (6.3) are
\(\bigoplus_{V\in\widehat G}\mathbb Z\) in degree zero and zero in
degree one. For example a projection over its unitization differs from
its scalar projection by a norm-small tail; when the tail norm is below
one, the standard near-projection unitary homotopy removes that tail.
Thus its K-class is supported in finitely many blocks. The same tail
removal applies to unitaries, and matrix stability gives the stated
groups. The finite-dimensional matrix computation has \(K_0=\mathbb Z\),
\(K_1=0\).

The cycle comparison of Theorem 4.3 sends the finite even representation
\(V\) to a rank-one projection in its multiplicity block in (6.3). Hence
the index \([\ker T]-[\ker T^*]\) is independent of all cycle relations
and distinct irreducibles give independent classes. This proves (6.2).

The zero-operator cycle on \(V\otimes W\) is a product of the two finite
cycles: its creation operators are compact and its positivity condition
for the zero first operator is zero. Thus multiplication is tensor
product. The ordinary flip \(V\otimes W\to W\otimes V\) is an
intertwining unitary, proving commutativity.

For an arbitrary equivariant cycle, tensoring by \(V\) with diagonal
action gives the scalar action \[
[V]\cdot x=[V]\boxtimes x.
\tag{6.4}
\] Associativity and bilinearity in Theorem 2.3 prove the module laws.
The signed flip moves an even \([V]\) through another factor without a
sign; naturality proves that all equivariant pullbacks, pushforwards and
products respect (6.4). \(\square\)

\textbf{Proposition 6.2 (the ring for a locally compact group).} For any
second countable locally compact \(G\), \(R^*(G)\) is a graded
commutative unital ring acting on the equivariant groups as graded
scalars. Its degree-zero part acts without parity signs.

\textbf{Proof.} Theorem 2.6 supplies composition, associativity and
identity \([\mathbb C]\). For scalar source and coefficient the external
and composition products agree by their exterior-extension definition.
The signed flip is the usual identification of scalar factors, so
symmetry reads \(xy=(-1)^{ij}yx\) for \(x\in R^i(G)\), \(y\in R^j(G)\).
Exterior product with an algebra identity gives its scalar action.
Associativity and interchange prove the module laws and compatibility
with products, and naturality proves compatibility with algebra maps.
This does not identify the ring with finite-dimensional representation
theory for a noncompact group. \(\square\)

\subsection{7. Restriction and compact
induction}\label{restriction-and-compact-induction}

A continuous homomorphism \(r:H\to G\) gives a restriction map \[
\begin{gathered}
r^*:KK_G^i(A,B)\\
\longrightarrow KK_H^i(A,B).
\end{gathered}
\tag{7.1}
\] by using the actions through \(r\) on the same cycle. Norm continuity
of \(h\mapsto r(h)\cdot F\), covariance and (1.4) follow by composition.
The same operation on interval modules proves well-definedness.
Composition of homomorphisms gives composition of restrictions. It
preserves products: the same operator remains a product cycle for the
restricted actions, so uniqueness in Theorem 2.6 applies. This includes
quotient homomorphisms, not only subgroup inclusions.

For a closed subgroup \(H\subset G\), let \(\operatorname{Ind}_H^G B\)
consist of continuous functions \(b:G\to B\) satisfying \[
\begin{gathered}
b(gh)=\beta_{h^{-1}}b(g),\\
\|b(g)\|\longrightarrow0\quad\text{on }G/H,\\
(k\cdot b)(g)=b(k^{-1}g).
\end{gathered}
\tag{7.2}
\] The inverse on \(h\) is part of the right-coset convention. When
\(B\) is a G-algebra restricted to \(H\), the map
\(b(g)\mapsto\beta_g b(g)\) identifies (7.2) with
\(C_0(G/H)\widehat\otimes B\) carrying the diagonal action. Its
covariance is checked by \(\beta_{gh}\beta_{h^{-1}}=\beta_g\).

For an equivariant Hilbert \(B\)-module \(E\), use the analogous section
condition \(\eta(gh)=U_h^{-1}\eta(g)\), with pointwise multiplication
and inner product. Denote the resulting module by
\(\operatorname{Ind}_H^G E\).

\textbf{Lemma 7.1 (induced sections for compact groups).} Suppose \(G\)
is compact and \(H\) is closed. Evaluation at each \(g\) has dense range
in \(E\), and induced compact operator fields are precisely \[
\mathcal K(\operatorname{Ind}_H^G E)
 \cong\operatorname{Ind}_H^G\mathcal K(E).
\tag{7.3}
\] The induced module is countably generated.

\textbf{Proof.} For \(\xi\in E\) and \(f\in C(G)\), put \[
\eta(g)=\int_H U_h\xi\,f(gh)\,dh.
\tag{7.4}
\] Changing the variable \(h=h_0^{-1}k\) shows
\(\eta(gh_0)=U_{h_0}^{-1}\eta(g)\). On \(H\), choose a nonnegative
continuous integral-one function supported near the identity. Extend it
continuously to \(G\), which is possible since compact metrizable \(G\)
is normal and \(H\) is closed. Equation (7.4) at the identity
approximates \(\xi\), by continuity of \(h\mapsto U_h\xi\). Left
translates do the same at any \(g\). This proves density of all
evaluation ranges and uses no local section of \(G\to G/H\).

A rank-one operator of two induced sections is the field
\(g\mapsto\theta_{\eta(g),\zeta(g)}\), which has the required covariance
and norm continuity. This embeds the left side of (7.3) into the right
side. To prove density, fix a continuous equivariant compact field
\(k(g)\). At one fibre approximate its value by a finite sum of rank-one
operators, and approximate the two vectors of each by evaluations of
global induced sections, as just proved. The error remains small on a
neighborhood, since its norm is continuous on \(G\) and descends to
\(G/H\). Cover compact \(G/H\) by finitely many such neighborhoods and
combine the approximations with a scalar partition of unity. A
nonnegative partition function can be split into its two square roots on
the two sections, so each patched term is still a rank-one operator.
This proves uniform density and hence (7.3).

Finally \(E\) is a separable Banach space under our standing hypotheses.
The induced module is a closed subspace of separable \(C(G,E)\), hence
separable. A countable norm-dense family generates it as a module, since
a coefficient approximate identity converges on every vector. Thus it is
countably generated. \(\square\)

\textbf{Theorem 7.2 (compact induction).} For compact \(G\) and closed
\(H\), induction gives \[
\begin{gathered}
\operatorname{Ind}_H^G:KK_H^i(A,B)\\
\longrightarrow KK_G^i(\operatorname{Ind}_H^G A,\\
       \operatorname{Ind}_H^G B).
\end{gathered}
\tag{7.5}
\] It preserves products, homomorphism classes and successive induction.

\textbf{Proof.} Average an input cycle to make \(F\) H-invariant. On its
induced module use the pointwise source action and pointwise operator
\(F\). The latter is well-defined because it commutes with \(H\), and is
invariant under left translation by \(G\). Each localized defect is the
continuous compact field obtained by applying the original defect to the
source section \(a(g)\). It is compact by (7.3). An invariant interval
homotopy induces the same construction over the interval coefficient
algebra, so this operation is well-defined and additive.

For products the pointwise tensor map \[
(\eta\widehat\otimes\zeta)(g)
 \longmapsto\eta(g)\widehat\otimes\zeta(g)
\tag{7.6}
\] identifies the tensor product of induced modules with the induced
tensor module. It is isometric by the pointwise inner product, and onto
by fibrewise density followed by the finite-neighborhood partition
argument of Lemma 7.1. An induced invariant product operator has the
connection conditions first for the pointwise creation sections, by
(7.3), then for every vector by density. Its localized positivity also
holds in the induced quotient: the pointwise positive quotient field has
a continuous positive lift modulo compact fields. To see this last
assertion directly, take the self-adjoint represented localized
anticommutator \(L(g)\). Its negative part is compact in each fibre
because its quotient is positive. Functional calculus makes \(L(g)_-\) a
continuous compact field with the required covariance; (7.3) makes it
compact on the induced module. Thus its operator quotient is positive.
Product uniqueness proves compatibility with (7.5).

Induction of a homomorphism cycle is the pointwise induced homomorphism
cycle, so identities are preserved. For \(L\subset H\subset G\),
successive sections are identified by \[
(\Psi\eta)(g)(h)=\eta(gh).
\tag{7.7}
\] The inverse evaluates the inner section at \(e\). The two covariance
conditions give the inverse identities; compactness of the quotients and
the supremum inner product give an isometric onto map. It intertwines
source actions, group actions and invariant operators, proving
successive induction. \(\square\)

\textbf{Theorem 7.3 (compact Frobenius reciprocity).} For compact \(G\),
closed \(H\), a G-algebra \(A\) and an H-algebra \(B\), there is a
natural isomorphism \[
\begin{gathered}
KK_H^i(A,B)\\
\cong KK_G^i(A,\operatorname{Ind}_H^G B).
\end{gathered}
\tag{7.8}
\] The forward map induces and represents \(a\in A\) at \(g\) by
\(\phi(\alpha_{g^{-1}}a)\). The inverse restricts to \(H\) and pushes
coefficients forward by evaluation at \(e\).

\textbf{Proof.} The described forward source action is covariant and
preserves the induced section condition. It is the pullback of (7.5) by
the equivariant map \(a\mapsto(g\mapsto\alpha_{g^{-1}}a)\) into
\(\operatorname{Ind}_H^G A\). Evaluation at \(e\) is H-equivariant,
because \(b(h^{-1})=\beta_h b(e)\). Thus both maps are defined on cycles
and on homotopies.

Their composite on an H-cycle is identified with the original by the
unitary \[
\begin{gathered}
\operatorname{Ind}_H^G E
 \widehat\otimes_{\operatorname{ev}_e}B
 \longrightarrow E,\\
\eta\widehat\otimes b\longmapsto\eta(e)b.
\end{gathered}
\tag{7.9}
\] Its inner product is exactly the evaluated inner product, and Lemma
7.1 gives dense range. It intertwines the representation, H-action and
invariant operator.

For the other composite take a G-cycle on \(Y\), with coefficient
\(D=\operatorname{Ind}_H^G B\), and average its operator invariantly.
Set \(E=Y\widehat\otimes_{\operatorname{ev}_e}B\), with the induced
H-action. Write \([y]_e\) for its fibre quotient. Define \[
(Wy)(g)=[U_g^{-1}y]_e.
\tag{7.10}
\] This is norm continuous, has the H-covariance condition, and obeys \[
\begin{gathered}
\langle Wy,Wz\rangle(g)
 =\langle y,z\rangle_D(g),\\
W(yb)(g)=(Wy)(g)b(g).
\end{gathered}
\tag{7.11}
\] Hence it is an isometry. Its fibre values are dense: at a fixed
\(g\), the quotient of \(U_g^{-1}Y\) is dense in \(E\). This quotient
statement can be expressed as the ordinary Hilbert-module fibre
quotient, so no identity in nonunital \(B\) is being inserted.
Finite-neighborhood approximation and a partition of unity on compact
\(G/H\) show that the range is dense in the induced module.
Multiplication of \(y\) by a partition function is permitted through the
central multiplier action of \(C(G/H)\) on \(D\). Therefore \(W\) is
onto.

For the group action, \(W(U_k y)(g)=Wy(k^{-1}g)\). For the source
action, \(W(\phi(a)y)(g)=\phi_e(\alpha_{g^{-1}}a)Wy(g)\). An invariant
\(F\) passes to \(F_e\) and satisfies \(W(Fy)(g)=F_e Wy(g)\). Thus \(W\)
identifies the second composite cycle with the original. These same
unitaries over interval modules prove the inverse identities on homotopy
classes. The auxiliary Clifford factor is untouched, giving both
degrees. \(\square\)

\subsubsection{Induction for a locally compact
group}\label{induction-for-a-locally-compact-group}

\textbf{Lemma 7.4 (local averaging and induced compact fields).} For any
closed \(H\subset G\), there is a nonnegative continuous function
\(\alpha\) on \(G\) such that \[
\int_H\alpha(gh)\,dh=1,
\tag{7.12}
\] and its support over each compact subset of \(G/H\) is compact in
\(G\). Induced sections have dense evaluation ranges, are countably
generated, and obey (7.3), now with compact fields vanishing at infinity
on \(G/H\).

\textbf{Proof.}
\href{supporting/noncompact-foundations/noncompact-foundations.html\#ncf-002}{NCF.2}
proves that the quotient is locally compact Hausdorff and second
countable.
\href{supporting/noncompact-foundations/noncompact-foundations.html\#ncf-001}{NCF.1}
gives its compactly supported partitions and cutoffs;
\href{supporting/noncompact-foundations/noncompact-foundations.html\#ncf-006}{NCF.6}
proves all continuity and support assertions in the following averaging
construction. Choose nonnegative \(f_j\in C_c(G)\) whose orbit integrals
\(F_j(gH)=\int_H f_j(gh)\,dh\) are positive on an open cover of \(G/H\).
Such a function is obtained at each coset by taking a bump positive at a
representative. Each orbit integral is continuous: near a fixed
representative, its integration variable lies in a common compact subset
of \(H\). The quotient map is open, so invariant continuity descends to
the quotient. Its support is in the compact set
\(q(\operatorname{supp}f_j)\).

Take a locally finite partition of unity \(\chi_j\) with supports
contained in the positivity sets, and put \[
\alpha(g)=\sum_j
 \chi_j(gH)\frac{f_j(g)}{F_j(gH)}.
\tag{7.13}
\] Each summand is extended by zero outside that positivity set. Local
finiteness gives continuity, (7.12) follows by integration, and a
compact quotient set meets only finitely many partition supports. Over
that set the support of \(\alpha\) therefore lies in finitely many
compact supports of \(f_j\). It is closed there and hence compact. In
particular \(g\mapsto\alpha(g\,\cdot)\) is locally continuous in
\(L^1(H)\), by uniform continuity on its common compact integration
support.

For evaluation density use (7.4) with compactly supported \(f\). By
\href{supporting/noncompact-foundations/noncompact-foundations.html\#ncf-004}{NCF.4},
nonnegative functions of integral one exist on \(H\) with arbitrarily
small compact support.
\href{supporting/noncompact-foundations/noncompact-foundations.html\#ncf-001}{NCF.1}
extends each such function from closed \(H\) to a compactly supported
continuous function on \(G\), retaining nonnegativity. The resulting
section is bounded, with norm supported on \(q(\operatorname{supp}f)\),
by the continuous orbit-integral bound. It is therefore induced, and
gives the same density proof as before. A compactly supported induced
compact field is approximated by the finite-neighborhood argument of
Lemma 7.1. Cutoffs on the quotient approximate every vanishing field,
proving (7.3).

Every compact quotient set has a compact set of representatives: cover
it by finitely many images of precompact open subsets of \(G\), and take
their compact closures. For sections supported in a fixed compact
quotient set, the supremum norm is thus detected on a compact subset of
\(G\). Restriction to \(L\) is an isometric injection of this section
space into \(C(L,E)\). The latter is separable by NCF.1: the compact
subspace \(L\) is second countable, and countable generation over
separable \(B\) makes \(E\) separable by approximating the finite linear
combinations of a countable generating family using a countable dense
subset of \(B\). A subspace of a separable metric space is separable:
balls of rational radius about a countable dense set give a countable
base; choose a point in each basic set meeting the subspace. A countable
quotient exhaustion and scalar cutoffs, NCF.1, therefore give a dense
union of separable section spaces. The induced module is separable. If
\(\eta_j\) is a countable norm-dense set of its vectors, the right
\(B\)-linear span of these vectors is dense: each
\(\eta_j e_n\to\eta_j\) for an approximate identity of \(B\), by the
inner-product norm calculation in the Hilbert-module foundations. Hence
it is countably generated. \(\square\)

\textbf{Theorem 7.5 (general induction).} The induction map (7.5), its
naturality, product compatibility and successive-induction formula hold
for every closed subgroup of a second countable locally compact group.

\textbf{Proof.} An H-cycle has a self-adjoint contractive H-continuous
operator \(F\). On its induced module first form the bounded pointwise
operator field \[
F_\alpha(g)=\int_H(h\cdot F)\alpha(gh)\,dh.
\tag{7.14}
\] The local \(L^1\) continuity in Lemma 7.4 makes this a
norm-continuous field. Changing \(h=h_0^{-1}k\) gives
\(F_\alpha(gh_0)=h_0^{-1}\cdot F_\alpha(g)\), so it acts adjointably on
induced sections. It is bounded by one. For each source value \(a(g)\),
its difference from \(F\) is two-sided locally compact for
\(\phi(a(g))\), by the H-equivariance defect and (7.12). On compact
subsets of the quotient these compact differences vary continuously,
using the common compact integration supports. The actual cycle defects
of \(F_\alpha\) are therefore continuous compact fields, and they vanish
at infinity after the source section acts, because \(\|a(g)\|\to0\). The
same is true of the localized difference between any fixed left
translate of \(F_\alpha\) and \(F_\alpha\). Uniform boundedness handles
the quotient tail, and the compact-support integral handles local
continuity.

This field need not have a norm-continuous orbit uniformly over the
whole quotient. Choose nonnegative \(f\in C_c(G)\) of integral one and
define \[
F_{\alpha,f}=\int_G f(k)(k\cdot F_\alpha)\,dk.
\tag{7.15}
\] The strict integral is adjointable. Its orbit is norm continuous by
the \(L^1\)-translation bound from (2.15). Its difference from
\(F_\alpha\) is locally compact for the induced source: the preceding
translated defects are norm-continuous compact module operators on the
compact integration support. Expanding square and commutator defects now
proves that (7.15) is an equivariant cycle. Two choices of \(\alpha,f\)
have a two-sided locally compact difference, since at each fibre they
are integral-one averages of H-translates of \(F\). Those compact fields
are continuous and vanish after multiplication by the source. Both
resulting operators are G-continuous, so the straight path is an
equivariant cycle. Thus the class is independent of the choices. The
same construction for an interval coefficient module proves homotopy
invariance and additivity. It commutes with coefficient maps and source
pullbacks, proving naturality.

For product compatibility use the tensor unitary (7.6), now by quotient
cutoffs and Lemma 7.4. Let \(F\) be an H-product of \(F_1,F_2\). A
translate \(h\cdot F\) is a connection for \(h\cdot F_2\), which differs
locally compactly from \(F_2\). Hence \(h\cdot F-F\) is a
zero-connection, as well as a two-sided locally compact source
perturbation. At each fibre the averages (7.14)--(7.15) therefore
preserve the connection conditions for the correspondingly averaged
\(F_2\). Creation errors are continuous compact fields and vanish at
quotient infinity by the vanishing norm of the creation section.
Averaging \(F_1\) changes its localized positivity tests by exactly the
locally compact input-perturbation calculation of Lemma 2.1 of the
product lesson. Thus the pointwise localized anticommutator is positive
modulo compact operators. Its negative part is a continuous compact
field vanishing at infinity; Lemma 7.4 makes that negative part compact
on the induced module. This is the required operator-quotient
positivity. The averaged \(F\) is consequently a G-product, and
uniqueness in Theorem 2.6 proves compatibility.

For successive induction the section unitary (7.7) retains its vanishing
conditions. A section vanishing on \(G/L\) restricts to a vanishing
section on each closed fibre \(H/L\), and the supremum over that fibre
vanishes on \(G/H\), since the image of a compact subset of \(G/L\) is
compact. Conversely, restrict a successively induced section to a
compact set of representatives for its relevant outer quotient support.
Its values there form a norm-compact family of inner sections, which has
a uniform vanishing tail on \(H/L\), by finite approximation. The
resulting compact representative set times the compact inner support
maps to a compact subset of \(G/L\). This proves vanishing of the
flattened section and the onto assertion of (7.7).

Under this unitary both averaged operators are integral-one averages of
translates of the original operator by the smallest subgroup. Their
difference is two-sided locally compact for the pointwise source. On
compact quotient sets this is a continuous compact field by the local
support property of Lemma 7.4; globally the source's vanishing tail
supplies the required decay. Both final operators have G-continuous
orbits, so the straight path identifies their classes. The zero operator
of a homomorphism cycle remains zero under each average, so induction
preserves homomorphism classes and identities. \(\square\)

\subsection{8. Computations}\label{computations}

The two-element kernel model in Section 1 now has its full cycle
interpretation. The next examples test three distinct features:
representation labels for a trivial action, loss of the stabilizer for a
free transitive action, and a grading changed by a reflection.

For finite \(G\) acting trivially on a trivially graded algebra \(B\),
\[
K_i^G(B)\cong R(G)\otimes_{\mathbb Z}K_i(B).
\tag{8.1}
\] Indeed \(B\rtimes G=B\otimes C^*(G)\), and the finite-dimensional
group algebra is a direct sum of matrix algebras, one for each
irreducible. This decomposition follows from the invariant orthogonal
decomposition and Schur's lemma in Proposition 6.1; the regular
representation is faithful by Proposition 4.1. Matrix stability gives
one copy of \(K_i(B)\) per irreducible, and (6.4) identifies its
representation-ring labels. There is no Künneth assumption on \(B\).

Let \(\mathbb T\) act on \(C(\mathbb T)\) by rotation. The homogeneous
space is \(\mathbb T/\{e\}\), so Theorem 7.3 gives \[
\begin{gathered}
K_0^{\mathbb T}(C(\mathbb T))=\mathbb Z,\\
K_1^{\mathbb T}(C(\mathbb T))=0.
\end{gathered}
\tag{8.2}
\] The same crossed-product groups follow by (4.8). The free proper
orbit module of \emph{Proper actions, free actions and the orbit space},
Proposition 9.2 and Theorem 9.5, identifies this crossed product with
\(\mathcal K(L^2(\mathbb T))\). A character of \(\mathbb T\) acts on the
integer group by its dimension, namely one: in Frobenius reciprocity it
restricts to the trivial subgroup.

For a Clifford example, let \(G=\mathbb Z/2\) act on \(C_1\) by sending
its odd self-adjoint generator \(e\) to \(-e\). The crossed product is
generated by \(e\) and an even \(u\), with \[
\begin{gathered}
e^2=u^2=1,\qquad ue=-eu,\\
e=\begin{pmatrix}0&1\\1&0\end{pmatrix},\\
u=\begin{pmatrix}1&0\\0&-1\end{pmatrix}.
\end{gathered}
\tag{8.3}
\] The four elements \(1,e,u,eu\) span the dense polynomial algebra and
their displayed matrices are linearly independent. Thus this is
\(M_2(\mathbb C)\), graded by conjugation with \(u\). The balanced
spinor \(\mathbb C^{1|1}\) gives its graded Morita equivalence with
\(\mathbb C\). Green--Julg therefore gives \[
\begin{gathered}
KK^G(\mathbb C,C_1)=\mathbb Z,\\
KK_G^1(\mathbb C,C_1)=0.
\end{gathered}
\tag{8.4}
\] The grading in (8.3) is needed for this conclusion. The
nonequivariant odd Clifford algebra has \(KK(\mathbb C,C_1)=0\); the
reflection action changes the equivariant group.

\subsection{9. Exercises}\label{exercises}

\textbf{9.1. Finite representations.} For a finite group \(G\), compute
\(KK^G(\mathbb C,\mathbb C)\) and its ring product. Give the result for
a finite cyclic group of order \(m\).

\textbf{9.2. A finite kernel proof.} Prove Green--Julg directly for a
finite group by computing the invariant matrix algebra on
\(\ell^2(G)\otimes B\). Include the normalization for Haar measure and
the cycle equivalence relations.

\textbf{9.3. The discrete dual.} Starting with an equivariant Fredholm
cycle for a countable discrete group, integrate it and recover its
action from the full crossed product. Prove that these operations
preserve homotopy, also for nonunital \(A\).

\textbf{9.4. A homogeneous space.} For compact \(G\) and closed \(H\),
prove \(K_0^G(C(G/H))\cong R(H)\) and compute the action of \(R(G)\).
Determine the odd group.

\subsection{10. Solutions}\label{solutions}

\textbf{Solution to 9.1.} Proposition 6.1 identifies every even cycle
with the virtual representation \([\ker T]-[\ker T^*]\) of its invariant
off-diagonal Fredholm operator. Distinct irreducibles are independent by
the compact-block decomposition (6.3) and the standard-cycle
isomorphism, so the group is the free abelian group on the irreducibles.
Tensoring two finite zero-operator cycles is their product, hence the
multiplication coefficients are the multiplicities in the tensor-product
decomposition.

For a cyclic group, diagonalize the unitary representing a generator.
Its eigenvalues are the \(m\)-th roots of unity, and its irreducibles
are the one-dimensional characters \(1,t,\ldots,t^{m-1}\). Tensoring
multiplies eigenvalues, so the ring is \[
R(\mathbb Z/m)=\mathbb Z[t]/(t^m-1).
\tag{10.1}
\] The corresponding odd group is zero by (6.3).

\textbf{Solution to 9.2.} Put \(n=|G|\) and give each point Haar mass
\(1/n\). A kernel operator is \((K\xi)(t)=n^{-1}\sum_r K(t,r)\xi(r)\).
The invariant condition from (4.6) gives \[
\begin{gathered}
K(t,r)=\beta_{t^{-1}}f(tr^{-1}),\\
f(s)=K(e,s^{-1}).
\end{gathered}
\tag{10.2}
\] Its composition is exactly the crossed-product convolution with
normalized counting measure, by substituting \(s=tr^{-1}\) in the finite
sum. Its adjoint is the crossed-product involution, by transposing and
adjointing this matrix kernel. Every entry lies in \(B\), so the
operator is in \(M_n(B)=\mathcal K(\ell^2(G)\otimes B)\); conversely
every invariant matrix has the form (10.2). Faithfulness follows either
by inspection of its \(e\)-row, followed by the finite-group full-norm
argument, or directly by embedding each covariant pair through its
constant regular vector as in Proposition 4.1. Thus the invariant matrix
algebra is the full crossed product.

For the same integrated operator, normalized Haar coefficients
\(f_{\mathrm{norm}}\) and counting-Haar coefficients
\(f_{\mathrm{count}}\) satisfy \[
\begin{gathered}
\frac1n\sum_g f_{\mathrm{norm}}(g)U_g
\\ =\sum_g f_{\mathrm{count}}(g)U_g,\\
f_{\mathrm{count}}=\frac1n f_{\mathrm{norm}}.
\end{gathered}
\tag{10.2a}
\] When \(B\) is unital, the identity coefficient is \(n1_B\delta_e\)
for normalized Haar and \(1_B\delta_e\) for counting Haar. This is
consistent with the normalized-Haar matrix identity kernel
\(K(t,r)=n\delta_{t,r}1_B\).

Add infinite trivial multiplicity, use Theorem 3.1 to standardize scalar
equivariant cycles, and use Lemma 4.2 to identify their invariant
adjointables with the multipliers of the invariant compact algebra.
Ordinary scalar cycles on the crossed product have the same multiplier
picture. Even unitaries, compact perturbations, norm operator paths and
degenerate operators match under this identification. Theorem 2.3 and
its non-equivariant counterpart identify these relations with module
homotopy. This proves the isomorphism on the actual K-groups in both
degrees; it is not just the displayed matrix-algebra identity.

\textbf{Solution to 9.3.} Define \(\Phi(\sum a_gu_g)=\sum\phi(a_g)U_g\).
The covariance identities give a representation of the finite-support
convolution algebra and its adjoint, and the full universal norm extends
it. Equation (5.2) and the localized equivariance defect give compact
commutators; the localized square and adjoint defects follow term by
term and then by norm approximation. Over an interval coefficient module
the same compact computations hold, so integration preserves homotopy.

For the inverse first pass to the essential representation using
Proposition 5.1 of the product lesson. Extend it to canonical multiplier
unitaries \(u_g\), and restrict to \(A\). Equation (5.2), applied to
both \(a\) and \(au_g\) in the crossed product, recovers the localized
equivariance defect. All group maps are continuous in the discrete
topology. For an interval homotopy take its essential replacement over
the interval first, then perform this same multiplier extension.
Endpoint replacements are locally compact perturbations of the essential
endpoint cycles, so they preserve the equivariant class by the straight
path described in Theorem 5.1. The two operations are inverse by the
integrated covariant correspondence. In the nonunital case \(u_g\) is in
the multiplier algebra but \(au_g\) is in the crossed product, and the
latter membership is exactly what was used in (5.2). No scalar identity
in \(A\) is required.

\textbf{Solution to 9.4.} Set \(A=B=\mathbb C\) in Theorem 7.3, with the
trivial H-action on \(B\). Equation (7.2) becomes
\(\operatorname{Ind}_H^G\mathbb C=C(G/H)\). Thus \[
K_i^G(C(G/H))\cong KK_H^i(\mathbb C,\mathbb C).
\tag{10.3}
\] Proposition 6.1 for compact \(H\) identifies these with \(R(H)\) and
zero. Explicitly an H-representation \(V\) gives the associated bundle
with sections \(\eta(gh)=\pi(h)^{-1}\eta(g)\); the induced zero-operator
cycle is its equivariant bundle class. Evaluation at the identity
recovers \(V\), so this description also identifies the isomorphism on
finite representatives. Tensoring the associated bundle with a
G-representation \(W\) changes its identity fibre to
\(\operatorname{Res}_H^G W\otimes V\). Hence the \(R(G)\)-action on
\(R(H)\) is restriction followed by multiplication.

\subsection{What this lesson does not
prove}\label{what-this-lesson-does-not-prove}

The non-equivariant connection, product, associativity,
essential-replacement and homotopy-comparison theorems used above are
proved in the exact preceding lessons named in the introduction. Full
crossed-product covariance and faithful regular-norm independence are
proved in the two named crossed-product prerequisites. Continuous
functional calculus, strict multiplier extension and the compact
spectral theorem are the basic analysis and Hilbert-module
prerequisites. The compact and noncompact Haar, cutoff, quotient and
integration inputs have the complete proofs linked in the introduction
and Lemma 7.4.

The two Green--Julg theorems have the compact and discrete hypotheses
stated above. General descent, its multiplicativity, proper actions and
assembly belong to the following lessons. Compact Frobenius reciprocity
(7.8) is not asserted here for arbitrary noncompact quotients.

The linear equivariant Bott equivalence is proved in Section 9 of
\emph{Thom isomorphisms and K-orientations in KK}, using the
compact-group product now constructed here. No Bott assertion was used
in the proof of Theorem 2.3.

\subsection{References}\label{references}

\begin{itemize}
\item
  B. Blackadar,
  \href{https://www.bruceblackadar.com/Mathematics/book6.pdf}{\emph{K-Theory
  for Operator Algebras}, author-posted corrected second edition}. The
  \href{https://www.bruceblackadar.com/mathpubs.html}{author identifies
  this freely downloadable version} as a slightly corrected second
  edition. The source retains its author copyright and the usage terms
  stated on that page. Sections 20.1--20.5, printed pp.~205--210, for
  equivariant cycles, stabilization, technical partitions, Green--Julg,
  representation classes and induction. Their additional equivariant
  proofs are supplied in Sections 1--7 above, including orbit
  continuity, the compact and discrete cycle comparisons, and local
  induction averages.
\item
  The names attached to these constructions are recorded in that freely
  accessible account: Kasparov for equivariant KK, Mingo and Phillips
  for equivariant stabilization, Julg for the compact-group isomorphism,
  and Wassermann for compact Frobenius reciprocity. The programme
  proofs, rather than references in that account, supply every result
  used here.
\end{itemize}

\end{document}
